\documentclass[11pt,a4paper]{article}
\usepackage[utf8]{inputenc}
\usepackage[T1]{fontenc}
\usepackage{lmodern,amsmath,amssymb,amsthm,mathtools,booktabs,enumitem,microtype}
\usepackage{natbib}
\usepackage[margin=1in]{geometry}
\usepackage{xr-hyper}
\usepackage{hyperref}
\externaldocument[IS-]{00_is}
\externaldocument[GE-]{01_law_of_graphic_equalisation_theory_referenced}
\externaldocument[PH-]{01a_graphic_equalisation_physics_grounding}
\externaldocument[NV-]{04_narrativium_theory_referenced}
\externaldocument[AD-]{09_adaptability}
\newtheorem{definition}{Definition}[section]
\newtheorem{theorem}[definition]{Theorem}
\newtheorem{proposition}[definition]{Proposition}
\newtheorem{corollary}[definition]{Corollary}
\theoremstyle{remark}\newtheorem{remark}[definition]{Remark}
\title{Graphic Equalisation --- Cosmology Grounding\\\large Horizons, Measures, and What Has No Outside}
\author{Matthew Parslow\\Independent Researcher}
\date{\textbf{Draft \texorpdfstring{\textemdash}{---}\ \PaperDraft}\\[2pt]\PaperDisplayDate \texorpdfstring{\textbullet}{\textperiodcentered}\ document revision \PaperRevision}
\usepackage{fancyhdr}
\pagestyle{fancy}
\fancyhf{}
\fancyfoot[C]{\thepage}
\fancyfoot[R]{\scriptsize\textsc{\PaperDraft\,r\PaperRevision}}
\renewcommand{\headrulewidth}{0pt}
\fancypagestyle{plain}{\fancyhf{}\fancyfoot[C]{\thepage}%
  \fancyfoot[R]{\scriptsize\textsc{\PaperDraft\,r\PaperRevision}}\renewcommand{\headrulewidth}{0pt}}
% GENERATED by tools/build_paper_metadata.py from paper_metadata.json.
% DO NOT EDIT: edit paper_metadata.json, then regenerate this file.
\providecommand{\PaperDraft}{draft-2026-09-06.1}
\providecommand{\PaperDisplayDate}{6 September 2026}
\expandafter\def\csname PaperRevision@00_is\endcsname{129}
\expandafter\def\csname PaperRevision@01_law_of_graphic_equalisation_theory_referenced\endcsname{294}
\expandafter\def\csname PaperRevision@01a_graphic_equalisation_physics_grounding\endcsname{47}
\expandafter\def\csname PaperRevision@01b_graphic_equalisation_mind_and_self_grounding\endcsname{47}
\expandafter\def\csname PaperRevision@01c_graphic_equalisation_cosmology_grounding\endcsname{11}
\expandafter\def\csname PaperRevision@01d_graphic_equalisation_born_recovery\endcsname{6}
\expandafter\def\csname PaperRevision@01e_graphic_equalisation_communication_grounding\endcsname{6}
\expandafter\def\csname PaperRevision@01f_graphic_equalisation_qualia_and_consciousness_grounding\endcsname{3}
\expandafter\def\csname PaperRevision@01g_graphic_equalisation_total_perspective_vortex_grounding\endcsname{4}
\expandafter\def\csname PaperRevision@01h_graphic_equalisation_information_gravity_quantum_reconciliation\endcsname{5}
\expandafter\def\csname PaperRevision@01i_graphic_equalisation_passive_continuation_momentum_gravity\endcsname{12}
\expandafter\def\csname PaperRevision@01s_graphic_equalisation_statement\endcsname{294}
\expandafter\def\csname PaperRevision@02_seldon_manifolds_theory_referenced\endcsname{31}
\expandafter\def\csname PaperRevision@02a_seldon_manifolds_psychological_grounding\endcsname{15}
\expandafter\def\csname PaperRevision@02b_seldon_manifolds_social_and_psychological_application\endcsname{7}
\expandafter\def\csname PaperRevision@02c_seldon_manifolds_decision_making\endcsname{11}
\expandafter\def\csname PaperRevision@03_law_of_narrative_momentum_theory_referenced\endcsname{32}
\expandafter\def\csname PaperRevision@04_narrativium_theory_referenced\endcsname{23}
\expandafter\def\csname PaperRevision@05_frames_relation_and_the_empty_limit\endcsname{11}
\expandafter\def\csname PaperRevision@06_assessing_an_account_of_this_kind\endcsname{29}
\expandafter\def\csname PaperRevision@07_causal_attitude_networks_on_the_substrate\endcsname{4}
\expandafter\def\csname PaperRevision@08_certified_answers_as_an_audit\endcsname{4}
\expandafter\def\csname PaperRevision@09_adaptability\endcsname{94}
\expandafter\def\csname PaperRevision@09a_adaptability_implementation\endcsname{7}
\expandafter\def\csname PaperRevision@10_applied_graphic_equalisation\endcsname{7}
\expandafter\def\csname PaperRevision@10a_red_queen_race_civilisational_rate_of_change_llms_biological_timescales_anthropocene_transition\endcsname{13}
\expandafter\def\csname PaperRevision@11_relational_information_ethics\endcsname{16}
\expandafter\def\csname PaperRevision@A1_constructive_graphic_equalisation_implementation_referenced\endcsname{8}
\expandafter\def\csname PaperGenerated@01s_graphic_equalisation_statement\endcsname{1}
\edef\PaperMetadataStem{\jobname}
\expandafter\ifx\csname PaperRevision@\PaperMetadataStem\endcsname\relax
  \PackageError{paper_metadata}{No metadata for \PaperMetadataStem}{Add the paper stem to paper_metadata.json.}
  \def\PaperRevision{?}
\else
  \expandafter\let\expandafter\PaperRevision\csname PaperRevision@\PaperMetadataStem\endcsname
\fi
\newcommand{\PaperDateLine}{\textbf{Draft \texorpdfstring{\textemdash}{---}\ \PaperDraft}\\[2pt]\PaperDisplayDate\ \texorpdfstring{\textbullet}{\textperiodcentered}\ document revision \PaperRevision}
\newcommand{\PaperStatementDateLine}{\textbf{Draft \texorpdfstring{\textemdash}{---}\ \PaperDraft}\\[2pt]\PaperDisplayDate\ \texorpdfstring{\textbullet}{\textperiodcentered}\ statement}
\expandafter\ifx\csname PaperGenerated@\PaperMetadataStem\endcsname\relax
  \newcommand{\PaperFooterStamp}{\textsc{\PaperDraft\,r\PaperRevision}}
\else
  \newcommand{\PaperFooterStamp}{\textsc{\PaperDraft\ / statement / from long r\PaperRevision}}
\fi
\newcommand{\PaperNoticeStamp}{\textbf{DRAFT --- \PaperDraft\ r\PaperRevision} (\PaperDisplayDate).}


\begin{document}
\maketitle
\begin{center}
\fbox{\parbox{0.92\textwidth}{\small
\textbf{DRAFT --- \PaperDraft r\PaperRevision} (\PaperDisplayDate).
This document is under active revision and is circulated for comment, not as a
finished result. It is a companion to \emph{The Law of Graphic Equalisation ---
Theory} (referred to below as Graphic Equalisation) and to its physics companion
\emph{The Law of Graphic Equalisation --- Physics Grounding} (the physics
companion), and inherits their definitions, epistemic boundary, and
level-of-claims. The other documents referred to below carry short names too, and
those are the forms used: \emph{Is} (the base paper), \emph{Narrativium}, and the
adaptability paper. Definitions, numbering and claims may change between drafts;
cite the version string above if referring to it. Sections marked as
falsifiability commitments are the intended points of attack.}}
\end{center}
\vspace{1em}

\begin{abstract}
This companion examines cosmological questions using frame-relative description.
It distinguishes the limits of observation from the topology and boundary
conditions of a cosmological model. A compact universe need not have an external
container. Horizon-limited data do not directly determine all unobserved regions;
inferring homogeneity from observed near-isotropy requires additional assumptions.
The measure problem in eternal inflation remains open. It does not follow that
all local, model-relative measures are ill-posed, or that a probability requires
an observer outside the universe. The paper offers conceptual distinctions and
conditional interpretations; it derives no cosmological parameter or empirical
prediction.
\end{abstract}

\section{Scope, and what this adds}\label{sec:scope}
The physics companion already holds cosmology's in-frame half, and the division between them is
principled rather than editorial. Questions asked \emph{from inside} a frame --- what the seed was, why
these constants, where the arrow of time comes from --- are physics questions, and that companion
interprets them conditionally: the Big Bang as a chosen entry rather than an intrinsic edge, the Past Hypothesis as a
seed and therefore rightly an assumption, fine-tuning inverted by self-location. None of it is
repeated here.

What is left is the boundary set, and this paper is short because that set is small: what the
subject of cosmology is, what a horizon is, what may be said of what lies past one, and what becomes of
probability when the sample space has no outside. The split follows the account's own line between
what is in a frame and what is between frames --- one companion for the questions with an answer
inside, one for the questions about the edge. A reader who finds this document thin has found that
division working rather than a gap in it, and the two are meant to be read in order.

Two disciplines carry over and are stated once. The account prices operations and does not
license them, and it distinguishes what is \emph{observed} from what is \emph{modelled} --- a
distinction \emph{Is} makes for its own object and which does most of the work below.

The result in one line, since the rest is the case for it. \emph{You cannot have the ensemble and
you cannot have a totality --- and this paper can say why.} That is the whole of what is on offer.
Cosmology already knows it cannot have them: the measure problem is called unsolved, homogeneity is
called an assumption, and both are called that by people with no interest in this account. Those
are correct instincts held without grounds, and the grounds are what is supplied here.

The difference is practical rather than decorative. An unsolved problem invites better
regularisations and more of them; an ill-posed one redirects the effort spent on that, which is
worth something to whoever would otherwise spend it. Knowing why you cannot have a thing is not the same
as being told again that you cannot.

\section{The universe is not an object}\label{sec:not-an-object}
Cosmology's ordinary subject does not exist, and \emph{Is} already says so. That paper gives
three readings which each yield nothing (\emph{Is}, Section~\ref{IS-sec:one-not-enough} and
Section~\ref{IS-sec:graph}): a single relatum with nothing to relate to, the empty case, and
\emph{everything at once read from outside, which has nothing left over to stand in relation to it}. The
universe as a whole is the third.

\begin{remark}[Totality is not an external observation frame]\label{prop:totality}
Under Definition~\ref{GE-def:frame}, an external tracking boundary requires
a specified relation. This prohibits treating a description of everything as
an external observation; it does not prohibit a mathematical model with
compact spatial sections and no spatial boundary, such as a closed FLRW
universe or a flat three-torus. Spatial topology and an epistemic frame boundary
are different objects.
\end{remark}
The claim below concerns observation and modelling, not an inference of
noncompact spatial topology. A whole compact cosmological model can be
well posed without a physical exterior. The base paper's coherence
proposition names two degenerate cases; the totality reading is discussed
in its Section~\ref{IS-sec:one-not-enough}.



The failure is not that the outside is empty but that there is always one more step outward, which is
the unbounded frontier of Graphic Equalisation's Section~\ref{GE-sec:scope}, read in the extensive
direction. This is a claim about \emph{frames}. What \emph{Is} posits --- everything, with no container, modelled
and not observed (its Section~\ref{IS-sec:one-way}) --- is not a frame and is not the object cosmology
individuates; the proposition says that no \emph{individuated} universe is the totality, which is the
observed/modelled distinction of Section~\ref{sec:scope} doing the work it was announced for.

The consequence is a correction and not a limitation. It is usual to say that we observe only
part of the universe and must infer the rest. On this account the observable universe is not a sample of
the object --- \emph{it is the object}, being what stands in relation to the frame doing the
cosmology. The horizon does not truncate a larger thing; it individuates the thing there is.

This is less revisionary than it sounds, and the practice already runs this way. The fits that
determine cosmological parameters are made to what is inside the horizon. A quantity that made no
difference to any fit was doing no work in the account, which is \emph{Is}'s own test.

\section{A horizon is a boundary relation}\label{sec:horizon}
The particle horizon bounds signals received since the initial epoch. The
event horizon bounds signals that can ever be received, and the Hubble radius
is a local expansion scale; they are not interchangeable. The discussion
here uses the particle horizon. Its growth can bring previously inaccessible
regions into view without a new local interaction at that region.

A particle horizon is the boundary relation itself, not a fence inside a larger space. Graphic
Equalisation's frame is an instantiated boundary $\mathcal{B}_{W:G}$ (Definition~\ref{GE-def:frame}),
and individuation is performed by the container rather than by the thing (Graphic Equalisation's
Remark~\ref{GE-rem:only-g}). A horizon individuates a cosmological frame in exactly that sense,
and its dependence on the observer's worldline is then expected rather than awkward: different
worldlines instantiate different boundaries, so they have different frames, and there is no frame
containing all of them from which to say which is right.

It does not follow that there is nothing outside, and the account says the opposite. A frame is
individuated by its container, so a frame that is individuated \emph{has} one --- \emph{the boundary
relation is itself a relation to the outside}, and a cosmological frame could not be a frame without it.
\emph{There is an outside. We are not in it.}

\begin{proposition}[The relation beyond a horizon runs up, not across]\label{prop:beyond}
What lies past a horizon is related to us \emph{through the superframe we are both subframes of} (the
physics companion's Corollary~\ref{PH-cor:nonlocal}); its content is a superframe dimension,
first-class where it lives (Graphic Equalisation's Proposition~\ref{GE-prop:higher-frame}) and
\textnormal{\textsc{null}} in our frame.
\end{proposition}

A connection between subframes runs up rather than across, and \textsc{null} here is \emph{not the
same status as absent}: the value is unreadable from here, but it was never a dimension of this frame,
so failing to read it discards nothing locally available. That distinction is the correction, and the
entangled case already makes it. A traced-out partner
is not information thrown away; the purity was never in the subframe. \emph{So what lies beyond a
horizon is real, related, and inaccessible} --- three properties ordinary usage runs together and the
account keeps apart.

What remains refused is only the reading. An extrapolation past the horizon is a model of a
superframe dimension made from below, which is a legitimate thing to hold and an illegitimate thing to
call measured. \emph{The prohibition is on the value, not on the existence}, and it costs the practice
nothing, since no measurement inside discriminates between competing values of a dimension that is
\textsc{null} here.

\begin{remark}[Where this could be wrong, and it is testable]\label{rem:bubble}
This is not idle, and on the corrected reading it is not awkward either. Searches for signatures
of bubble collisions in the microwave background \citep{feeney2011,planck2016} are searches for a
coupling to something past the horizon. A detection would refute nothing here. Measuring projects a
superframe dimension into the subframe's aperture, which is the physics companion's standing
description of what measuring does (its Corollary~\ref{PH-cor:nonlocal}) --- so a found signature is
a dimension becoming accessible and not an outside coming into being. The
aperture widens; the dimension was there. What it would refute is this paper's assumption about the
present aperture, which is an assumption about us and not about the world.
\end{remark}

\section{The cosmological principle divides}\label{sec:cosmological-principle}
Isotropy is observed. Homogeneity is inferred, and the inference has a modal leg. What is measured
is that the sky looks the same in every direction to high precision. What is wanted is that the universe
is the same everywhere, and the step between them is the Copernican principle: that our position is not
special among positions. \emph{Our position is not special} compares this position against positions
we do not occupy, and \emph{Is}'s commitment (its Section~\ref{IS-sec:graph}) is that what bears on a
claim is a demonstration and not an imagining, since imagining a world costs nothing and settles
nothing --- so what follows is not the claim that nothing is out there.

\begin{remark}[Isotropy does not alone establish homogeneity]\label{prop:copernican}
Isotropy along our past light cone alone does not entail homogeneity on a
spatial slice or beyond the particle horizon. Homogeneity is inferred through
models and constrained by light-cone consistency tests
\citep{clarkson2012,maartens2011}. Theory-mediated predictions are permitted;
this boundary concerns treating an untested extrapolation as a direct reading.
\end{remark}
The horizon problem is the observed near-isotropy of the CMB across patches
without a shared causal past at recombination in the unextended model.
Inflation supplies earlier causal contact and predicts large-scale homogeneity
alongside testable in-horizon consequences. Such an extrapolation is a
conditional physical prediction, not a forbidden frame-free measurement.



The standard derivation makes the point for us, inside its own hypothesis. The
Ehlers--Geren--Sachs theorem \citep{ehlers1968} obtains Friedmann--Lema\^{i}tre--Robertson--Walker
geometry not from isotropy measured here but from isotropy measured \emph{by a family of observers
everywhere}; its realistic form \citep{stoeger1995} weakens exactness and not universality, and the
observational programme built on it is explicit that what is being combined with the microwave
measurements is the Copernican principle \citep{clarkson2012,maartens2011}. So the modal step is
not appended to the derivation. It is the antecedent. The theorem is sound and the inference is valid;
what it requires is a measurement over a domain no frame reaches, which is the one thing that cannot be
supplied from inside.

This is the sharpest form of the section's claim and it is cosmology's own statement of it. The
route from isotropy to homogeneity is not sloppy and is not naive --- \emph{it is exact, and it is
conditional on a quantifier}, and the account's whole contribution here is to say why that quantifier
has no referent.

The field's own vocabulary already concedes this. Homogeneity is called an assumption, not a
result, and it is called that by people with no interest in this account. What is added here is
not the suspicion but the reason: it is not an assumption because the measurement is hard, but because
the quantity it asserts stands in no relation to the frame asserting it.

\section{The measure problem is the no-prior-free proposition}\label{sec:measure}
This is the case where the account does real work, and it is the case cosmology already knows is
broken.

Eternal inflation produces unboundedly many pocket universes \citep{guth2007}. To extract a prediction
--- what a typical observer sees, how likely a given constant is --- one needs a probability measure over
them. The theory does not supply one. Different regularisations give different answers, some give
absurd ones, and the difficulty is named and open \citep{freivogel2011}: the youngness paradox and the
Boltzmann-brain problem \citep{bousso2007} are what the disagreement looks like when a stipulation is
pushed to its consequences.

\begin{remark}[No frame-free measure is supplied by the counting vocabulary]\label{prop:measure}
A probability measure requires a declared ensemble and counting or weighting
rule. A standpoint with no relation to any member supplies none. This excludes
an unexplained frame-free denominator; it does not establish that a local
causal-patch measure is impossible or that existing regularisations cannot
converge. Constructed observer-relative measures remain legitimate candidates
whose physical assumptions must be compared \citep{bousso2007,freivogel2011}.
\end{remark}

The adaptability companion's Proposition~\ref{AD-prop:no-prior-free}
requires declared baselines; it does not deny that a uniform or local measure
can be constructed. Local causal-patch descriptions address the global
counting problem rather than exemplifying its impossibility. The convergence
of local and global constructions is consequently relevant evidence, not a
case the frame discipline excludes.

\begin{remark}[What this does not do to inflation]\label{rem:inflation}
Nothing here bears on inflation as physics. Inflationary models make horizon-scale predictions
that are tested inside a frame and stand or fall on those. What is refused is a specific inferential
move --- probability over an ensemble with no outside --- and a theory can be entirely correct while one
inference drawn from it is malformed. Conflating the two would be the error this paper is
otherwise about.
\end{remark}

\section{Anthropic reasoning inherits the same defect}\label{sec:anthropic}
Typicality arguments are the measure problem wearing a different hat, and worse than that: the
term they turn on names two different objects, and the argument survives only by sliding between them.

Read one way, \emph{universe} means the ensemble --- all of them, the population a typical
member is typical of. Read the other, it means what the word ordinarily means: the observable,
relative universe, the one a frame stands in relation to. Cosmological usage takes the second and infers
with the first.

\begin{proposition}[Typicality equivocates, and each reading fails separately]\label{prop:typical}
On the ensemble reading there is no population to be typical of (Remark~\ref{prop:measure}); on
the relative reading there is exactly one universe --- the frame's --- and \emph{typicality is
undefined over a set a thing exhausts}. \emph{Neither reading supports the inference.}
\end{proposition}

They fail for unrelated reasons --- the ensemble reading wants the counting standpoint
Remark~\ref{prop:measure} refuses, the relative reading a class its one member exhausts --- which
is why moving between them looks like an answer.

This is stronger than the usual complaint, which is that the reference class is hard to choose.
The difficulty is not that the class is hard to pick but that on one reading there is no class and on
the other the class has one member.

\emph{Is}'s care about \emph{everything} is the same care and should be repeated exactly.
It does not say that everything \emph{is} nothing --- it says the unqualified identity is asserted from a
standpoint at which it does not hold, and that what is true is a claim about \emph{readings}: everything
at once, read from outside, has nothing left over to stand in relation to it and so yields nothing.
So \emph{everything} does not name a population whose members could be counted. It names an
operation that returns nothing, and a distribution cannot be defined over the output of that operation.

Between the two zeros is where the distinctions live, which is \emph{Is}'s own sentence and
is the whole of what cosmology has to work with: not the ensemble, which reads as nothing, and not a
totality, which does not close --- but this frame, which is where everything measurable is.

The physics companion reaches the same place from the other direction and the two should be read
together: there, fine-tuning is dissolved by self-location
(Proposition~\ref{PH-prop:self-location}), since a different seed is a different
observer and there is no trans-seed standpoint from which to compare. Here the same absence removes the
distribution. One prohibition, met twice.

\section{Two measurements that disagree}\label{sec:tension}
A word on the Hubble tension, and it is a word of caution rather than a result. Early-universe and
late-universe determinations of the expansion rate disagree beyond their stated errors, and the
disagreement is the field's own open problem \citep{divalentino2021}.

The account says only what it says. By the adaptability paper's result a measurement is a
property of a four-place arrangement (Remark~\ref{AD-rem:chollet}), and two determinations made under
different arrangements are two
measurements rather than one measurement made twice --- so comparing them requires an argument that the
moved coordinate does not bear. That is exactly what the field is doing, at length and carefully,
which is why this section claims nothing further. The account does not dissolve the tension, does
not favour either value, and does not predict which coordinate is at fault. It says the comparison is
the work, which cosmologists established without help.

\section{What is not claimed}\label{sec:not-claimed}
No cosmological parameter is derived here, and none could be. No observation is predicted. Nothing
above distinguishes between inflationary models, settles the expansion tension, or bears on dark energy's
equation of state.

Nor is any of the negative work a criticism of practice. Homogeneity is already called an
assumption; the measure problem is already called unsolved; the tension is already treated as a
comparison problem. The account agrees with the field's instincts and supplies reasons for them,
which is worth less than a prediction and more than nothing.

\section{Falsifiability commitments}\label{sec:falsifiability}

\noindent\emph{On the tags.} Each commitment below is marked \textsc{[shape]} or \textsc{[detail]}.
A \textsc{[detail]} failing is a revision: the account survives with that mechanism replaced. A
\textsc{[shape]} failing costs \emph{range} rather than a mechanism, because what fails is something
the substrate is claimed to require or to permit --- so the account does not hold where it claimed to,
which is the more serious of the two and is still a boundary rather than an annihilation.

\begin{enumerate}[leftmargin=*]
\item \textsc{[detail]}~\textbf{An aperture wider than assumed.} A confirmed signature of a bubble
collision, or of any structure causally beyond the particle horizon, would show a superframe dimension
to be accessible from here. That refutes neither the account nor this paper's reading of what lies
outside --- both hold that it is there --- but it refutes this paper's assumption about the
\emph{present aperture}, which is the load-bearing assumption of Section~\ref{sec:horizon} and is about
us rather than about the world. What changes on refutation is stateable: that section's claim that no
measurement inside discriminates between competing values of a dimension that is
\textnormal{\textsc{null}} here would be false for that dimension, and the prohibition on reading its
value would lift for that dimension alone. Marked \textsc{[detail]} because what is replaced is an
assumption about the present width of our aperture, not a claim about the substrate.

\item \textsc{[detail]}~\textbf{Model-relative measures.} A proposed measure must specify its
ensemble, weighting and regulator or local sampling rule, and test robustness
and observational predictions. Agreement among measures would be progress on
the physical problem, not a refutation of the requirement to specify premises.

\item \textsc{[detail]}~\textbf{Anisotropy that survives.} Isotropy is the leg this paper keeps. A
large-scale anisotropy that survives scrutiny would remove the observed half of the cosmological
principle. It does not touch Remark~\ref{prop:copernican}, which declines the inference from
isotropy to homogeneity and is left standing \emph{a fortiori} if the antecedent fails; what it costs is
this paper's sentence that isotropy is measured, and with it the surviving-leg framing of
Section~\ref{sec:cosmological-principle}. Marked \textsc{[detail]} for that reason, and listed rather
than dropped because that framing is what the section is built on.

\item \textsc{[detail]}~\textbf{Typicality on a single reading.} Proposition~\ref{prop:typical} fails if
a typicality argument can be stated on one reading of \emph{universe} held throughout --- either a
countable ensemble carrying a frame-internal measure, which is item~2 again, or a relative universe
whose reference class has more members than the one the frame exhausts --- without sliding between the
two. The claim is the equivocation, so exhibiting the argument without it is what refutes it.

\item \textsc{[shape]}~\textbf{Derivation of a seed} (inherited). If a cosmological initial condition
is ever derived rather than installed --- from within, with no boundary condition imported --- then the
physics companion's seed result (Proposition~\ref{PH-prop:past-hyp}) and this paper's dependence on it
both fail. The commitment belongs to that companion and is carried here because this paper leans on it.

\item \textsc{[detail]}~\textbf{Totality and topology.} A claim about the universe's physical
extent requires a cosmological model and topology. The absence of an external
epistemic vantage does not prohibit a compact spatial model without a boundary;
Remark~\ref{prop:totality} makes that distinction explicit.

\end{enumerate}

\noindent The horizon claim is epistemic. Correlations in a joint quantum state
(Corollary~\ref{PH-cor:nonlocal}) supply no theorem about cosmological regions
beyond an observer's horizon. A superframe description there is a proposed
model, not an inherited Bell result. Likewise, homogeneity can be tested within
specified model assumptions without reading an otherwise inaccessible dimension.

\begin{thebibliography}{99}

\bibitem[Bousso and Freivogel(2007)]{bousso2007}
Raphael Bousso and Ben Freivogel.
\newblock A paradox in the global description of the multiverse.
\newblock \emph{Journal of High Energy Physics}, 2007(06):018, 2007.

\bibitem[Clarkson(2012)]{clarkson2012}
Chris Clarkson.
\newblock Establishing homogeneity of the universe in the shadow of dark energy.
\newblock \emph{Comptes Rendus Physique}, 13:682--718, 2012.

\bibitem[Di Valentino et~al.(2021)]{divalentino2021}
Eleonora Di~Valentino, Olga Mena, Supriya Pan, Luca Visinelli, Weiqiang Yang,
  Alessandro Melchiorri, David~F. Mota, Adam~G. Riess, and Joseph Silk.
\newblock In the realm of the Hubble tension --- a review of solutions.
\newblock \emph{Classical and Quantum Gravity}, 38(15):153001, 2021.

\bibitem[Ehlers et~al.(1968)]{ehlers1968}
J\"{u}rgen Ehlers, P. Geren, and Rainer~K. Sachs.
\newblock Isotropic solutions of the Einstein--Liouville equations.
\newblock \emph{Journal of Mathematical Physics}, 9(9):1344--1349, 1968.

\bibitem[Feeney et~al.(2011)]{feeney2011}
Stephen~M. Feeney, Matthew~C. Johnson, Daniel~J. Mortlock, and Hiranya~V. Peiris.
\newblock First observational tests of eternal inflation.
\newblock \emph{Physical Review Letters}, 107:071301, 2011.

\bibitem[Freivogel(2011)]{freivogel2011}
Ben Freivogel.
\newblock Making predictions in the multiverse.
\newblock \emph{Classical and Quantum Gravity}, 28:204007, 2011.

\bibitem[Guth(2007)]{guth2007}
Alan~H. Guth.
\newblock Eternal inflation and its implications.
\newblock \emph{Journal of Physics A}, 40:6811--6826, 2007.

\bibitem[Maartens(2011)]{maartens2011}
Roy Maartens.
\newblock Is the universe homogeneous?
\newblock \emph{Philosophical Transactions of the Royal Society A}, 369:5115--5137, 2011.

\bibitem[Planck Collaboration(2016)]{planck2016}
Planck Collaboration.
\newblock Planck 2015 results XVI: Isotropy and statistics of the CMB.
\newblock \emph{Astronomy and Astrophysics}, 594:A16, 2016.

\bibitem[Stoeger et~al.(1995)]{stoeger1995}
William~R. Stoeger, Roy Maartens, and George F.~R. Ellis.
\newblock Proving almost-homogeneity of the universe: an almost Ehlers--Geren--Sachs theorem.
\newblock \emph{Astrophysical Journal}, 443:1--5, 1995.

\end{thebibliography}

\end{document}
